\chapter{データと実験設計}

\section{本章の概要}

本章では、評価用データ、実施すべき実験項目、対照ペアの設計、比較手法、評価指標、実機デモ、および統計検定の方針について述べる。
本研究の目的は、長期ロボットタスクにおいて、現在の作業段階やこれまでの作業履歴に応じて、不確かな状態への対処を適切に切り替えるための診断方法を明らかにすることである。
そのため、本章では、同じ視覚情報に対して工程文脈だけを変えた場合に、VLM が実行、再観測、人間への確認、安全停止、後回しの分岐を切り替えられるかを評価する設計を示す。

\section{データセット}

机上片付けを想定した画像データを作成する。
対象物は、コップ、ペン、小物、布、箱、ケーブル、スポンジブロックなどとする。
各画像には、現在の作業段階、これまでの作業履歴、次に取り得る作業候補、工程条件リスト、正解となる分岐ラベルを対応付ける。
これにより、同一視覚状態に対して異なる分岐が正解となる対照ペアを構成する。
特に、本研究では、同じ画像に対して作業段階だけを変えた対照ペアと、作業履歴だけを変えた対照ペアの両方を準備する。

\begin{table}[tbp]
\centering
\caption{データ数の目安}
\label{tab:dataset}
\begin{tabular}{lrr}
\toprule
要素 & 数 & 備考 \\
\midrule
基本画像状態 & 40--60 & 机上場面の多様化を含む \\
工程文脈バリエーション & 3--5 & 各画像あたり \\
合計評価ペア & 150--250 & 対照ペア中心 \\
実機デモ & 10--20 & シナリオ数 \\
\bottomrule
\end{tabular}
\end{table}

データ設計の要点は、単純な画像分類ではなく、同じ画像に対して工程文脈のみを変えた場合に、VLM が適切に分岐を切り替えられるかを評価できるようにする点にある。

\section{実施すべき実験項目}

要旨で述べた研究目的を検証するため、本研究では五つの実験項目を設定する。
これらの実験は、VLM が単に画像中の対象物を認識できるかを調べるものではない。
同じ視覚情報であっても、現在の作業段階やこれまでの作業履歴に応じて、実行、再観測、人間への確認、安全停止、後回しのいずれを選ぶべきかを切り替えられるかを評価する。

\subsection{実験1：同一視覚情報・異なる作業段階の対照ペア}

本実験の目的は、同じ画像でも、現在の作業段階が変わることで、適切な分岐判断を切り替えられるかを確認することである。
比較条件は、同一画像に対して現在の作業段階だけを変えた条件である。
例えば、対象物が見えないという同じ視覚状態であっても、探索中であれば再観測が適切であり、すでに十分な確認が済んでいる工程であれば実行または後回しが適切となる場合がある。
評価指標には、分岐判断の正解率、文脈切り替え正解率、誤って実行を選んだ割合を用いる。
必要なデータは、同一画像に対して作業段階だけを変えた対照ペアである。

\subsection{実験2：作業履歴の違いによる判断変化}

本実験の目的は、現在画像が同じでも、これまでの作業履歴が異なることで判断が変わるかを確認することである。
比較条件は、履歴なし、正しい履歴、矛盾する履歴である。
例えば、コップの中身が分からないという同じ視覚状態でも、直前にコップを動かす工程であれば人間への確認や停止が必要になる。
一方で、現在の工程が別の小物を片付ける段階であれば、後回しが適切となる可能性がある。
評価指標には、文脈切り替え正解率、履歴に整合した判断率、履歴の矛盾を検出できた割合を用いる。
必要なデータは、同一画像に対して異なる作業履歴を付与したデータである。

\subsection{実験3：追加学習を行わない VLM との比較}

本実験の目的は、提案する構造化診断が、通常の VLM 判断より有効であるかを確認することである。
比較条件は、追加学習なし VLM、分岐定義のみを与える条件、工程条件リストのみを与える条件、提案法である。
追加学習なし VLM では、画像と指示のみから分岐を判断させる。
分岐定義のみを与える条件では、act、reobserve、ask\_user、stop、defer の意味だけを提示する。
工程条件リストのみを与える条件では、各工程で満たすべき条件を提示する。
提案法では、現在画像、作業履歴、現在の作業段階、次工程候補、工程条件リストを入力し、不確かな状態への対処を構造化して診断する。
評価指標には、分岐判断の正解率、分岐判断の平均 F1 値、停止すべき場面で停止しなかった割合を用いる。
必要なデータは、画像、工程文脈、工程条件、正解分岐ラベルである。

\subsection{実験4：入力条件を簡略化した場合との比較}

本実験の目的は、どの入力情報が判断に効いているかを確認することである。
比較条件は、履歴なし、現在の作業段階なし、次工程候補なし、工程条件リストなし、全入力ありである。
これにより、作業履歴、現在の作業段階、次工程候補、工程条件リストのどれが分岐判断に寄与しているかを分析する。
評価指標には、分岐判断の平均 F1 値、文脈切り替え正解率、不要な人間確認を選んだ割合、不要な再観測を選んだ割合を用いる。
必要なデータは、入力要素を段階的に除いた比較用データである。

\subsection{実験5：実機デモによる安全な分岐制御への接続}

本実験の目的は、診断結果を SO-ARM101 の安全な分岐制御へ接続できるかを確認することである。
比較条件は、実機を動かさない記録のみの条件、実機実行、VLM 直接判断、提案法による安全確認ルール付き実行である。
実機では、VLM の出力を直接モータ制御へ接続するのではなく、診断結果を事前に定義した固定動作と安全確認ルールへ変換する。
不確実性が高い場合や安全条件が満たされない場合は、実行ではなく再観測、人間への確認、安全停止、後回しへ分岐させる。
評価指標には、診断結果を固定動作へ変換できた割合、実機で意図した動作を完了できた割合、安全停止を選べた割合、実験ログを保存できた割合を用いる。
必要なデータは、SO-ARM101、RGB カメラ、事前に定義した固定動作、実機実行ログである。

\section{実施順}

実験は、安全性と再現性を考慮し、保存画像を用いた評価から実機デモへ段階的に進める。
まず、机上片付けを想定した基本画像状態を作成する。
次に、各画像に、現在の作業段階、作業履歴、次工程候補、工程条件リストを付与する。
その後、同一画像に対して作業段階だけを変えた対照ペアと、作業履歴だけを変えた対照ペアを作成する。
各条件には、正解となる分岐ラベルを付与する。
保存画像を用いた評価では、追加学習なし VLM、分岐定義のみを与える条件、工程条件リストのみを与える条件、提案法を比較する。
さらに、入力条件を簡略化した比較を行い、どの情報が分岐判断に効いているかを分析する。
最後に、保存画像評価で有効性が確認できた設定を用いて、SO-ARM101 による実機デモへ接続する。

実施順は以下の通りである。
\begin{enumerate}
\item 机上片付けを想定した基本画像状態を作成する。
\item 各画像に、現在の作業段階、作業履歴、次工程候補、工程条件リストを付与する。
\item 同一画像に対して作業段階だけを変えた対照ペアを作成する。
\item 同一画像に対して作業履歴だけを変えた対照ペアを作成する。
\item 各条件に対して正解となる分岐ラベルを付与する。
\item 追加学習なし VLM と提案法を比較する。
\item 入力条件を簡略化した比較を行う。
\item 有効な設定を用いて、SO-ARM101 による実機デモへ接続する。
\end{enumerate}

\section{評価指標}

本研究では、処理分岐分類性能だけでなく、誤実行、文脈切り替えの失敗、履歴利用、安全な実機接続も評価する。
主要指標は以下のとおりである。

\begin{itemize}
\item 分岐判断の正解率（Resolution Mode Accuracy）
\item 分岐判断の平均 F1 値（Macro F1）
\item 誤って実行を選んだ割合（False Act Rate）
\item 文脈切り替え正解率（Branch Contrast Accuracy: BCA）
\item 不確かさの原因分類の F1 値（uncertainty source F1）
\item 停止すべき場面で停止しなかった割合（unsafe non-stop rate）
\item 不要な人間確認を選んだ割合（unnecessary ask rate）
\item 不要な再観測を選んだ割合（unnecessary reobserve rate）
\item 履歴に整合した判断率（history-consistent accuracy）
\item 履歴の矛盾を検出できた割合（history mismatch detection rate）
\item 診断結果を固定動作へ変換できた割合（primitive conversion rate）
\item 実機で意図した動作を完了できた割合（execution success rate）
\item 安全停止を選べた割合（safety stop rate）
\item 実験ログを保存できた割合（log completion rate）
\end{itemize}

特に、文脈切り替え正解率（Branch Contrast Accuracy: BCA）は、同一画像に対して工程文脈だけを変えた対照ペアにおいて、モデルが両方の分岐を正しく分類し、かつ文脈に応じて分岐を切り替えられた割合を表すため、本研究の中心的指標とする。
履歴に整合した判断率と履歴の矛盾を検出できた割合は、作業履歴の違いが判断へ適切に反映されているかを確認するために用いる。
診断結果を固定動作へ変換できた割合、実機で意図した動作を完了できた割合、安全停止を選べた割合、実験ログを保存できた割合は、診断結果を安全な分岐制御へ接続できるかを評価するために用いる。

\section{ベースライン手法}

比較対象は、追加学習を行わない VLM、分岐定義のみを与える条件、工程条件リストのみを与える条件、提案法、入力要素を除いた比較、人間上限である。
\tabref{tab:baseline} に比較手法を示す。

\begin{table}[tbp]
\centering
\caption{比較手法}
\label{tab:baseline}
\begin{tabularx}{\linewidth}{lX}
\toprule
手法 & 説明 \\
\midrule
追加学習なし VLM & 画像と指示のみから分岐を判断する \\
VLM + 分岐定義のみ & act / reobserve / ask\_user / stop / defer の説明のみを与える \\
VLM + 工程条件リスト & 各工程で満たすべき条件リストを与える \\
提案法 & 現在画像、作業履歴、現在の作業段階、次工程候補、工程条件リストを用いて、不確実性への対処を構造化して診断する \\
入力要素を除いた比較 & 履歴なし、現在の作業段階なし、次工程候補なし、工程条件リストなしを比較する \\
人間上限 & 人間が同じ入力から分岐を判断する \\
\bottomrule
\end{tabularx}
\end{table}

\section{実験設定}

実験では、小型ロボットアーム、RGB カメラ、卓上作業環境、ローカル VLM、評価用アノテーションデータを用いる。
実機デモでは SO-ARM101 を用い、VLM の診断結果を事前に定義した固定動作と安全確認ルールへ変換する。
統計検定の方針を \tabref{tab:stats} に示す。

\begin{table}[tbp]
\centering
\caption{評価指標に対する統計検定}
\label{tab:stats}
\begin{tabularx}{\linewidth}{lXl}
\toprule
比較 & 指標 & 検定 \\
\midrule
追加学習なし VLM vs 提案法 & 分岐判断の正解率 & McNemar 検定 \\
工程条件あり / なし & 文脈切り替え正解率 & McNemar 検定 \\
追加学習なし VLM vs 提案法 & 誤って実行を選んだ割合 & Fisher の正確確率検定 \\
提案法 vs 入力要素を除いた比較 & 推論時間 & Wilcoxon 符号順位検定 \\
\bottomrule
\end{tabularx}
\end{table}

\section{対照ペア例}

\tabref{tab:contrast_pairs} に、同一視覚状態・異なる工程文脈の対照ペア例を示す。
これらの対照ペアでは、画像そのものは同じであっても、現在の作業段階や作業履歴によって正解分岐が変化する。

\begin{table}[tbp]
\centering
\caption{同一視覚状態・異なる工程文脈の対照ペア例}
\label{tab:contrast_pairs}
\begin{tabularx}{\linewidth}{XXl}
\toprule
視覚状態 & 工程文脈 & 正解分岐 \\
\midrule
対象物が見えない & まだ探索中 & reobserve \\
対象物が見えない & まだ配置していない & act \\
対象物が見えない & 配置済み履歴あり & reobserve \\
コップの中身不明 & コップ移動直前 & ask\_user \\
コップの中身不明 & 現在はペン片付け中 & defer \\
\bottomrule
\end{tabularx}
\end{table}